\subsection{Local cohomology over a Macaulay ring}

In this section, we assume that A is a local Macaulay ring; we set \(\dim(A) = d\) .

The ideals generated by a sequence of elements of \(\mathfrak{m}_{\mathrm{A}}\) completely secant for A and of length \(d = \dim (\mathrm{A})\) form a cofinal subset \(\mathcal{D}_{cs}\) in the set \(\mathcal{D}\) of defining ideals of A. Indeed, let \((x_{1},\ldots ,x_{d})\) be a sequence of elements of \(\mathfrak{m}_{\mathrm{A}}\) completely secant for A (§ 2, n° 3, prop. 3); for every integer \(n\) , the sequence \((x_1^n,\dots ,x_d^n)\) is completely secant for A (A, X, p.158, prop. 6, c)), and generates a defining ideal (VIII, § 3, n° 2, cor. of prop. 3 and th. 1) contained in \(\mathfrak{m}_{\mathrm{A}}^n\) .

Let \(\mathfrak{a} \in \mathcal{D}_{cs}\) , and let \(\pi: L \to A/\mathfrak{a}\) be a free resolution of finite type, zero in degree \(>d\) (for example the Koszul complex associated with a completely secant sequence for A generating \(\mathfrak{a}\) ). Consider the dual \(L^{*} = \operatorname{Homgr}_{A}(L,A)\) of \(L\) ; since the depth of A is equal to \(d\) , one has \(\operatorname{Ext}_{A}^{i}(A/\mathfrak{a}, A) = 0\) for \(i < d\) ( \(\S\) 1, \(n^{\circ}\) 1, cor. 2 of prop. 2). Since \(L^{*}\) is of length \(\leqslant d\) , we deduce that \(H^{i}(L^{*})\) is zero for \(i \neq d\) and that \(H^{d}(L^{*})\) is identified with \(\operatorname{Ext}_{A}^{d}(A/\mathfrak{a}, A)\) (A, X, p.100, th. 1). Consequently, we have a homomorphism

\[
\pi^ {*}: \mathrm{L} ^ {*} (- d) \longrightarrow \operatorname{Ext} _ {\mathrm{A}} ^ {d} (\mathrm{A} / \mathfrak {a}, \mathrm{A})
\]

which defines a free resolution of finite type of \(\operatorname{Ext}_{\mathrm{A}}^{d}(\mathrm{A}/\mathfrak{a},\mathrm{A})\) .

Let M be an A-module; consider the canonical isomorphisms (loc. cit.)

\[
\varphi (\mathrm{L}, \mathrm{M}): \mathrm{H} (\operatorname{Homgr} _ {\mathrm{A}} (\mathrm{L}, \mathrm{M})) \rightarrow \operatorname{Ext} _ {\mathrm{A}} (\mathrm{A} / \mathfrak {a}, \mathrm{M})
\]

\[
\psi (\mathrm{M}, \mathrm{L} ^ {*} (- d)): \mathrm{Tor} ^ {\mathrm{A}} (\mathrm{M}, \mathrm{Ext} _ {\mathrm{A}} ^ {d} (\mathrm{A} / \mathfrak {a}, \mathrm{A})) \to \mathrm{H} (\mathrm{M} \otimes_ {\mathrm{A}} \mathrm{L} ^ {*}) (- d).
\]

Since the complex L is free of finite type, the canonical morphism of complexes \(M \otimes_{A} L^{*} \to \text{Homgr}_{A}(L, M)\) is an isomorphism; from this we deduce an isomorphism of graded A-modules \(H(M \otimes_{A} L^{*}) \to H(\text{Homgr}_{A}(L, M))\) . By composing the preceding isomorphisms, we obtain an isomorphism of graded A-modules, said to be canonical.

\[
\tau (\mathrm{L}, \mathrm{M}): \operatorname{Tor} ^ {\mathrm{A}} (\mathrm{M}, \operatorname{Ext} _ {\mathrm{A}} ^ {d} (\mathrm{A} / \mathfrak {a}, \mathrm{A})) (d) \longrightarrow \operatorname{Ext} _ {\mathrm{A}} (\mathrm{A} / \mathfrak {a}, \mathrm{M})
\]

which induces, for each integer i, an isomorphism

\[
\tau^ {i} (\mathrm{L}, \mathrm{M}): \operatorname{Tor} _ {d - i} ^ {\mathrm{A}} (\mathrm{M}, \operatorname{Ext} _ {\mathrm{A}} ^ {d} (\mathrm{A} / \mathfrak {a}, \mathrm{A})) \longrightarrow \operatorname{Ext} _ {\mathrm{A}} ^ {i} (\mathrm{A} / \mathfrak {a}, \mathrm{M}).
\]

For \(M = A\) , \(\tau^d(L, A)\) is the canonical isomorphism from \(A \otimes_A \operatorname{Ext}_A^d(\mathrm{A} / \mathfrak{a}, A)\) onto \(\operatorname{Ext}_\mathrm{A}^d(A / \mathfrak{a}, A)\) .

Let \(\mathfrak{b}\) be an ideal of \(\mathcal{D}_{cs}\) contained in \(\mathfrak{a}\) . Let \(\rho: R \to A/\mathfrak{b}\) be a finite free resolution of length \(\leqslant d\) and let \(p_{\mathfrak{ab}}: A/\mathfrak{b} \to A/\mathfrak{a}\) be the canonical surjection. By A, X, p.49, prop. 3, there exists a morphism of complexes \(\mathrm{P}_{\mathrm{LR}}: R \to L\) such that \(\pi \circ \mathrm{P}_{\mathrm{LR}} = p_{\mathfrak{ab}} \circ \rho\) . By prop. 2 of A, X, p.103, we have a commutative diagram

\[
\begin{array}{c c c} \operatorname{Tor} ^ {A} (M, \operatorname{Ext} _ {A} ^ {d} (A / \mathfrak {a}, A)) (d) & \xrightarrow {\operatorname{Tor} (1 _ {M} , \operatorname{Ext} ^ {d} (p _ {a b} , 1 _ {\mathrm{A}}))} & \operatorname{Tor} ^ {\mathrm{A}} (M, \operatorname{Ext} _ {A} ^ {d} (A / \mathfrak {b}, A)) (d) \\ \psi (L ^ {*} (- d), M) \Bigg \downarrow & & \Bigg \downarrow \psi (R ^ {*} (- d), M) \\ H (M \otimes_ {A} L ^ {*}) & \xrightarrow {H (1 _ {M} \otimes^ {t} P _ {L R})} & H (M \otimes_ {A} R ^ {*}) \\ \Bigg \downarrow & & \Bigg \downarrow \\ H (\operatorname{Homgr} _ {\mathrm{A}} (L, M)) & \xrightarrow {H (\operatorname{Homgr} (P _ {L R} , M))} & H (\operatorname{Homgr} _ {A} (R, M)) \\ \varphi (L, M) \Bigg \downarrow & & \Bigg \downarrow \varphi (R, M) \\ \operatorname{Ext} _ {\mathrm{A}} (A / \mathfrak {a}, M) & \xrightarrow {\operatorname{Ext} (p _ {a b} , 1 _ {M})} & \operatorname{Ext} _ {A} (A / \mathfrak {b}, M) \end{array}
\]

It follows first, by taking \(\mathfrak{a} = \mathfrak{b}\) , that the isomorphism \(\tau(\mathrm{L},\mathrm{M})\) does not depend on the choice of the resolution L of A/ \(\mathfrak{a}\) ; let us denote it \(\tau_{\mathfrak{a}}(\mathrm{M})\) . It then follows that the \(\tau_{\mathfrak{a}}(\mathrm{M})\) for \(\mathfrak{a} \in \mathcal{D}_{cs}\) form an inductive system of isomorphisms. Passing to the inductive limit, one obtains for each integer \(i\) , taking into account A, X, p.70, prop. 8, an isomorphism of A-modules

\[
\tau^ {i} (\mathrm{M}): \operatorname{Tor} _ {d - i} ^ {\mathrm{A}} (\mathrm{M}, \mathrm{H} _ {\mathrm{A}} ^ {d} (\mathrm{A})) \longrightarrow \mathrm{H} _ {\mathrm{A}} ^ {i} (\mathrm{M}).
\]

For M = A, \(\tau^{d}(A)\) is the canonical isomorphism from \(A \otimes_{A} H_{A}^{d}(A)\) onto \(H_{A}^{d}(A)\) .

Remarks.--- 1) Let \(f: \mathrm{M} \to \mathrm{N}\) be a homomorphism of A-modules. Using A, X, p.103, prop. 2, one proves that the following diagrams are commutative:

\[
\begin{array}{c c c} \operatorname{Tor} _ {d - i} ^ {\mathrm{A}} (\mathrm{M}, \mathrm{H} _ {\mathrm{A}} ^ {d} (\mathrm{A})) & \xrightarrow {\tau^ {i} (\mathrm{M})} & \mathrm{H} _ {\mathrm{A}} ^ {i} (\mathrm{M}) \\ \operatorname{Tor} _ {d - i} (f, 1) \Bigg \downarrow & & \Bigg \downarrow \mathrm{H} ^ {i} (f) \\ \operatorname{Tor} _ {d - i} ^ {\mathrm{A}} (\mathrm{N}, \mathrm{H} _ {\mathrm{A}} ^ {d} (\mathrm{A})) & \xrightarrow {\tau^ {i} (\mathrm{N})} & \mathrm{H} _ {\mathrm{A}} ^ {i} (\mathrm{N}) \quad . \end{array}
\]

\begin{enumerate}
\def\labelenumi{\arabic{enumi})}
\setcounter{enumi}{1}
\tightlist
\item
  Let
\end{enumerate}

\[
0 \to \mathrm{M} \to \mathrm{N} \to \mathrm{P} \to 0\tag{ℰ}
\]

an exact sequence of A-modules. Using A, X, p.104, prop. 3 and p.106, prop. 4, one proves that the following diagrams are commutative:

\[
\begin{array}{c c c} \operatorname{Tor} _ {d - i} ^ {\mathrm{A}} (\mathrm{P}, \mathrm{H} _ {\mathrm{A}} ^ {d} (\mathrm{A})) & \xrightarrow {\tau^ {i} (\mathrm{P})} & \mathrm{H} _ {\mathrm{A}} ^ {i} (\mathrm{P}) \\ \partial_ {d - i} (\mathscr {E}, \mathrm{H} _ {\mathrm{A}} ^ {d} (\mathrm{A})) \Bigg \downarrow & & \Bigg \downarrow \partial^ {i} (\mathscr {E}) \\ \operatorname{Tor} _ {d - i - 1} ^ {\mathrm{A}} (\mathrm{M}, \mathrm{H} _ {\mathrm{A}} ^ {d} (\mathrm{A})) & \xrightarrow {\tau^ {i + 1} (\mathrm{M})} & \mathrm{H} _ {\mathrm{A}} ^ {i + 1} (\mathrm{M}) \end{array} .
\]

\begin{enumerate}
\def\labelenumi{\arabic{enumi})}
\setcounter{enumi}{2}
\tightlist
\item
  Consider the isomorphism
\end{enumerate}

\[
\tau^ {d} (\mathrm{M}): \mathrm{M} \otimes_ {\mathrm{A}} \mathrm{H} _ {\mathrm{A}} ^ {d} (\mathrm{A}) \longrightarrow \mathrm{H} _ {\mathrm{A}} ^ {d} (\mathrm{M}).
\]

For \(x \in \mathrm{M}\) let us denote by \(f_x\) the map \(a \mapsto ax\) from A into M. For every \(u \in \mathrm{H}_{\mathrm{A}}^{d}(\mathrm{A})\) , one has

\[
\tau^ {d} (\mathrm{M}) (x \otimes u) = \mathrm{H} _ {\mathrm{A}} ^ {d} (f _ {x}) (u).
\]

This indeed follows from Remark 1 applied to the homomorphism \(f_{x}:A\to M\) .

